\documentclass[11pt,a4paper]{article}

\usepackage[margin=1in]{geometry}
\usepackage{setspace}
\usepackage{graphicx}
\usepackage{booktabs}
\usepackage{caption}
\usepackage{subcaption}
\usepackage{amsmath, amssymb}
\usepackage{enumitem}
\usepackage{hyperref}
\usepackage{multirow}
\usepackage{array}

\hypersetup{
  colorlinks=true,
  linkcolor=blue,
  citecolor=blue,
  urlcolor=blue
}

\title{Project Proposal: AI-based Damage Detection for Power Grid Equipment with SaaS-oriented Design}
\author{Student: Yujie Feng \quad Supervisor: Jin Hou\\
SWJTU-LEEDS Joint School, Computer Science}
\date{\today}

\begin{document}
\maketitle
\onehalfspacing

\section*{Administrative Information}
\begin{itemize}[leftmargin=*]
  \item \textbf{Student Name}: Yujie Feng
  \item \textbf{Area of Interest}: Artificial Intelligence, Image Detection, Deep Learning, Large-scale Model
  \item \textbf{Title}: AI-based Damage Detection for Power Grid Equipment with SaaS-oriented Design
  \item \textbf{Supervisor}: Jin Hou
  \item \textbf{School of Computing Resources}: Cloud servers optimized for AI/GenAI/HPC
  \item \textbf{Other Resources Needed}: Cloud server
  \item \textbf{Relevant Modules}: [Machine Learning], [Computer Vision], [Deep Learning], [Data Mining / Big Data Systems], [Distributed Systems] % replace as appropriate
\end{itemize}

\section{Motivation and Problem Statement}
Critical power grid assets (e.g., transmission lines, towers, substations) require effective monitoring to prevent outages and safety incidents. This project targets a \emph{SaaS-oriented AI framework} for \textbf{damage detection and monitoring}, capable of handling multimodal inputs (e.g., RGB/thermal images, sensor telemetry) and outdoor noise factors (lighting/weather/clutter). The scope includes instance-level detection/segmentation (e.g., nests/foreign objects, cracks, corrosion) and a cloud-ready backend to support multi-tenant operations.

\section{Objectives}
\begin{enumerate}[leftmargin=*]
  \item Design and prototype an \textbf{instance segmentation} pipeline (Mask R-CNN) for power grid anomalies (initial case study: bird nests on lines/towers).
  \item Achieve strong validation performance (\textbf{bbox} AP@\([0.50:0.95]\) $\ge$ 0.70, \textbf{segm} AP@\([0.50:0.95]\) $\ge$ 0.70 on internal validation).
  \item Build a \textbf{SaaS-oriented backend} prototype (auth, RBAC, tenant isolation; Python inference service via REST).
  \item Provide \textbf{reproducible evaluation} (COCO metrics, PR curves, IoU distributions) and presentation-ready visuals.
\end{enumerate}

\section{Related Work}
Mask R-CNN-based instance segmentation is widely used for infrastructure inspection under complex scenes. Cloud-native (SaaS) designs enable scalable deployments, tenant isolation, and lifecycle management of AI models. Domain adaptation and robust augmentation address outdoor domain shifts.

\section{Approach and Methodology}
\subsection{Data \& Annotations}
\begin{itemize}[leftmargin=*]
  \item COCO-format dataset (\verb|Datasets/my_coco_dataset/dataset.json|); positive images include polygonal masks (e.g., bird nests), \textbf{negative samples} are added to \texttt{images} without annotations.
  \item Extremely high-resolution sources (up to \(8688\times5792\)); training inputs resized to \(512\) (optionally \(768/1024\) when resources permit).
\end{itemize}

\subsection{Model and Training}
\begin{itemize}[leftmargin=*]
  \item \textbf{Model}: \texttt{maskrcnn\_resnet50\_fpn\_v2} with pretrained backbone; heads replaced for custom classes.
  \item \textbf{Augmentations}: Color jitter, grayscale, equalize, posterize, horizontal flip; \texttt{CustomRandomIoUCrop} guarded by a \emph{SafeIoUCrop} wrapper to skip cropping on negative samples.
  \item \textbf{Optimization}: AdamW (lr \(5\cdot10^{-4}\)), OneCycleLR, 40 epochs (early-stop ready).
  \item \textbf{Logging}: Per-epoch history to \texttt{history.jsonl}; best checkpoint metadata in \texttt{training\_metadata.json}.
\end{itemize}

\subsection{SaaS-oriented Design (Prototype)}
\begin{itemize}[leftmargin=*]
  \item \textbf{Architecture}: Python microservice (Flask/FastAPI) for inference; backend with RBAC, tenant isolation; REST APIs; Docker for containerization (Kubernetes-ready).
  \item \textbf{Security/Isolation}: JWT-based auth, tenant-level data separation, secure model registry endpoints.
\end{itemize}

\section{Experimental Setup}
\begin{itemize}[leftmargin=*]
  \item \textbf{Hardware}: NVIDIA GeForce RTX 3060 Laptop GPU (6\,GB VRAM), AMD Ryzen mobile CPU, 16\,GB RAM, NVMe SSD, Windows 11. During training the GPU memory stabilized around 4.2/6\,GB, CPU utilization peaked at ~40\%.
  \item \textbf{Software}: PyTorch 2.x, TorchVision 0.16 (with AMP and \texttt{cudnn.benchmark}), pycocotools, Jupyter, and custom COCO utilities for data loading and visualization.
  \item \textbf{Dataset}: 54 COCO images (43 train / 11 validation). All samples are patrol photographs of bird nests (positives) and empty power lines (negatives) from the same acquisition batch.
  \item \textbf{Training Config}: Input size 512$\times$512, batch size 4, AdamW ($\text{lr}=5\cdot 10^{-4}$), OneCycleLR for 40 epochs. Random IoU crop (trials=50), color/grayscale/equalize/posterize/horizontal-flip augmentations. On Windows the script is refactored into a \texttt{main()} entry point; DataLoader parallelism controlled via \texttt{NUM\_WORKERS} (default 2) with \texttt{persistent\_workers} and \texttt{prefetch\_factor}=2.
\end{itemize}

\section{Preliminary Results}
Unless otherwise noted, the following statistics are from the run directory \texttt{pytorch-mask-r-cnn-bird-nest/2025-11-12\_17-37-45}.

\subsection{COCO Metrics (Validation)}
\begin{center}
\begin{tabular}{lcc}
\toprule
Metric & \textbf{bbox} & \textbf{segm} \\
\midrule
AP@[0.50:0.95] & \(\mathbf{0.805}\) & \(\mathbf{0.795}\) \\
AP50 & \(1.000\) & \(1.000\) \\
AP75 & \(1.000\) & \(0.919\) \\
AR@[0.50:0.95] & \(0.837\) & \(0.828\) \\
\bottomrule
\end{tabular}
\end{center}
\textit{Interpretation}: High AP50/AP75 indicate accurate localization and mask quality; strong AR suggests low miss rate on the validation set.

\subsection{Training Dynamics and Observations}
\begin{itemize}[leftmargin=*]
\item \textbf{Loss Curve}: \texttt{train\_loss} drops from 1.41 to 0.14 while \texttt{valid\_loss} decreases from 0.65 to 0.19; both plateau around epoch 20 (see Fig.~\ref{fig:loss}). Early stopping could be triggered near \texttt{valid\_loss}$\approx 0.22$.
\item \textbf{GPU Utilization Improvements}: Initially the single-thread DataLoader caused oscillating GPU load (0\%/99\%). After refactoring into \texttt{main()}, setting \texttt{NUM\_WORKERS=2}, and enabling \texttt{persistent\_workers}/\texttt{pin\_memory}, the GPU stays at 70--85\% utilization.
\item \textbf{Memory \& I/O}: Reducing IoU crop trials from 100 to 50, combined with cuDNN autotuning and SSD storage, yields $\sim$50--60 seconds per epoch.
\item \textbf{Overfitting Risk}: Validation distribution mirrors the training domain; the model nearly replicates annotations (AP50=1.0, most IoUs in 0.9--1.0). Broader generalization requires new lines/seasons/backgrounds.
\end{itemize}

\subsection{Qualitative Examples}
\begin{figure}[h]
  \centering
  \includegraphics[width=\textwidth]{visualizations/image_1.png}
  \caption{GT (left, green) vs Prediction (right, red). Additional samples: see \texttt{visualizations/}.}
\end{figure}

\subsection{Online Inference Spot Checks}
\begin{itemize}[leftmargin=*]
  \item After deploying the trained model to the Flask inference service, we sampled several patrol photos that were never seen during validation; the detector consistently segmented bird nests and aligned bounding boxes with the actual locations, indicating encouraging generalization.
  \item Incoming production images span a wide range of resolutions (\(\leq 21\,\mathrm{MP}\)). We rescale them before inference and restore the coordinates proportionally during post-processing; once this fix is applied, front-end overlays match the offline evaluation view.
  \item Differences in season, lighting, or background between deployment scenarios and the curated dataset remain a concern. We plan to expand the corpus—especially hard negatives and heavy occlusion cases—to further harden robustness.
\end{itemize}

\subsection{Loss Curves and Performance Profiles}
\begin{figure}[h]
  \centering
  \includegraphics[width=0.95\textwidth]{visualizations/loss_curves.png}
  % \caption{Train/Validation loss for run 2025-11-12\_17-37-45. Early epochs plummeted 20 epochs after which the loss stabilized。}\label{fig:loss}
\end{figure}

\noindent\textbf{Score Threshold Sweeps:} Fig.~\ref{fig:threshold} shows precision/recall/F1 vs confidence threshold (IoU$\ge$0.5). Precision stays near 1.0 within 0.4--0.6; an operating threshold of 0.4--0.5 balances recall and false positives.

\subsection{Precision-Recall and IoU Distributions}
\begin{figure}[h]
  \centering
  \includegraphics[width=0.95\textwidth]{visualizations/pr_curve_bbox_ap50.png}
  \caption{PR Curve at IoU=0.50 for bounding boxes. Nearly horizontal precision near 1.0 indicates very few false positives as recall increases.}
\end{figure}

\begin{figure}[h]
  \centering
  \includegraphics[width=0.95\textwidth]{visualizations/pr_curve_segm_ap50.png}
  \caption{PR Curve at IoU=0.50 for segmentation masks.}
\end{figure}

\begin{figure}[h]
  \centering
  \includegraphics[width=0.95\textwidth]{visualizations/iou_hist_bbox.png}
  \caption{Bounding-box IoU Histogram. Right-skewed mass indicates strong localization (high overlaps).}
\end{figure}

\begin{figure}[h]
  \centering
  \includegraphics[width=0.95\textwidth]{visualizations/iou_hist_mask.png}
  \caption{Mask IoU Histogram (segm). High-IoU mass indicates accurate instance masks on validation samples.}
\end{figure}

\begin{figure}[h]
  \centering
  \includegraphics[width=0.95\textwidth]{visualizations/threshold_sweep_bbox.png}
  \caption{Score Threshold vs Precision/Recall/F1 for bbox at IoU$\ge$0.5. Operating points 0.4--0.5 offer the best trade-off.}\label{fig:threshold}
\end{figure}

\paragraph{Note on Figures.}
Replace placeholders with exports directly from \texttt{bird\_nest\_report.ipynb}. The notebook already computes COCO metrics, PR curves (AP50 slice), and IoU histograms (bbox + mask).

\section{Planned System Design and Deliverables}
\begin{itemize}[leftmargin=*]
  \item \textbf{Data/Model}: Expanded dataset with harder negatives; training pipeline with reproducible scripts; best checkpoint artifact and color map.
  \item \textbf{Service Layer}: Python inference microservice (REST), model versioning, RBAC, tenant isolation.
  \item \textbf{Evaluation Report}: Jupyter notebook with metrics, curves, histograms, and curated examples; PDF export for presentation.
\end{itemize}

\section{Risk Assessment and Mitigation}
\begin{itemize}[leftmargin=*]
  \item \textbf{Data scarcity / distribution shift}: Mitigate via negative samples, targeted augmentation, sliding-window inference on ultra-high-res inputs.
  \item \textbf{Overfitting}: Early stopping, weight decay, validation monitoring.
  \item \textbf{Scalability}: Containerized inference; plan for Kubernetes autoscaling if load increases.
  \item \textbf{Security}: RBAC + tenant isolation; secure token-based access; audit logs.
\end{itemize}

\section{Project Plan (Milestones)}
\begin{itemize}[leftmargin=*]
  \item Weeks 1--2: Dataset curation, COCO conversion, negative sample integration.
  \item Weeks 3--4: Baseline Mask R-CNN finetuning; initial evaluation notebook.
  \item Weeks 5--6: Hard example mining; model refinement; visualization polishing.
  \item Weeks 7--8: SaaS backend prototype (auth/RBAC/tenant isolation); inference microservice.
  \item Weeks 9--10: Integration, stress tests; final report and presentation.
\end{itemize}

\section{Conclusion}
Preliminary experiments show strong instance segmentation performance on the validation split (bbox AP@\([0.50:0.95]\) \(0.805\), segm \(0.795\); AP50~1.0), with robust qualitative results. The next steps are to expand and diversify data (especially negatives and hard cases), finalize SaaS-oriented components for multi-tenant deployment, and harden the evaluation with additional scenarios.

\bibliographystyle{plain}
\begin{thebibliography}{9}
\bibitem{he2017maskrcnn}
Kaiming He, et al. “Mask R-CNN.” ICCV, 2017.

\bibitem{lin2014coco}
Tsung-Yi Lin, et al. “Microsoft COCO: Common Objects in Context.” ECCV, 2014.

\bibitem{ren2015faster}
Shaoqing Ren, et al. “Faster R-CNN: Towards Real-Time Object Detection with Region Proposal Networks.” NeurIPS, 2015.
\end{thebibliography}

\end{document}